\scanpage{131}
\mspage{5.}
% archive sheet number: 130
\diagram{\begin{tikzcd}[ampersand replacement=\&,column sep=tiny]
X \arrow[r,"i_1(f)"] \arrow[dr,"f"'] \& X_1(f) \arrow[r,"i_1(p_1(f))"] \arrow[d,"p_1(f)"] \& X_2(f) \arrow[r] \arrow[dl,Rightarrow,"p_2(f)"] \& \cdots \& X_\alpha(f) \arrow[r,"i_1(p_\alpha(f))"] \arrow[dlll,"p_\alpha(f)"] \& X_{\alpha+1}(f) \arrow[r] \arrow[dllll,"p_{\alpha+1}(f)=p_1(p_\alpha(f))"] \& \cdots \\
\& Y \& \& \& \& \&
\end{tikzcd}}
\[\tag*{(1.14.)}\]
Donc, pour achever de définir la suite ordinale des $X_\alpha(f)$ \inserted{(munis \deleted{\unsure{de}} des $i_\alpha(f),p_\alpha(f)$)}, ou des $p_\alpha(f)$ (munis des $i_\alpha(f)$), il reste à décrire $X_1(f),i_1(f),p_1(f)$.

\deleted{\unsure{Je peux faire un}} La propriété essentielle de ce triple dont nous avons besoin est la suivante :

\deleted{Fact}

(1.15) $\operatorname{Fact}(C_0)$ (propriété de factorisation \unsure{p.r.|p.v.} à $C_0$)

Pour tout diagramme commutatif
\diagram{\begin{tikzcd}[ampersand replacement=\&,execute at end picture={
\node[anchor=east] at ([xshift=-5pt]current bounding box.west) {$D:$};}]
\deleted{X}A \arrow[r,Rightarrow,"\unsure{\xi|?|g}"] \arrow[d,"i"'] \& \deleted{B}X \arrow[d,"f"] \\
\deleted{X}B \arrow[r] \& Y
\end{tikzcd}}
avec $i\in C_0$

d'où un diagramme traits pleins
\deleted{il existe une factorisation}
\diagram{\begin{tikzcd}[ampersand replacement=\&,column sep=large]
A \arrow[r] \arrow[d,"i"'] \& \deleted{X} \arrow[r,Rightarrow] \& X \arrow[d,"i_D"] \arrow[dr,dashed,"i_1(f)"] \arrow[dd,bend right=28,dashed,"f"'] \& \\
B \arrow[rr] \& \& X_D\overset{\mathrm{def}}=\deleted{B}X\amalg_A B \arrow[r,dotted,"u_D"] \arrow[d,"p_D"] \& \deleted{\unsure{\Sigma_1|Z_1}}\unsure{X_1|X_\alpha}(f) \arrow[dl,dashed,"p_1(f)"] \\
\& \& Y \&
\end{tikzcd}}
, il existe un hom de la factorisation $f=p_Di_D$ de $f$ dans la factorisation $f=p_1(f)i_1(f)$, )
\deleted{$X$ \unsure{composé|complet}}
\deleted{i.e. une flèche dans $M_{/Y}$}

Pour \unsure{formuler} cette condition, j'ai utilisé l'existence des sommes \deleted{$\unsure{X_D|\Sigma_D|Z_D}$} \unsure{algamées|amalgamées} $\deleted{\unsure{\Sigma_D|Z_D}}X_D$, i.e. la condition \unsure{sur} $C_0$
